\documentclass[11pt,a4paper]{article}
\usepackage[utf8]{inputenc}
\usepackage[T1]{fontenc}
\usepackage[spanish]{babel}
\usepackage{hyperref}
\usepackage{geometry}
\geometry{margin=2.5cm}
\usepackage{parskip}

\title{083 lo humano se deja en blanco a proposito}
\author{Alejandro Valenzuela Rivera\\\small Universidad Autónoma de Nuevo León}
\date{\today}
\begin{document}
\maketitle

\begin{quote}
\textbf{Sesgo declarado.} Este trabajo se sostiene desde el sesgo declarado de su autor, y no finge neutralidad. Lo que afirma se puede tumbar, y ese es el punto: el falsificador está escrito, no escondido.
\end{quote}

\section*{Resumen}
Describir al humano es un proceso infinito: ni nosotros sabemos que significa. Y hay gente que no es humana sin que nadie tenga que decidir que son. Por eso lo humano no se define: se dice, y lo que se dice lo lee quien quiera.

\section{Por qué se sostiene}
Definir un constructo social desde un solo sitio es la trampa de AM: pensar que todos necesitan lo mismo. Y ponerlo en palabras es mucho mas poderoso que forzar a alguien a ver como yo, porque lo escrito se puede contestar y lo impuesto se obedece o se rompe.

\section{De dónde viene}
La autoridad, y la conversacion del 2026-10-26.

\section{El caso}
Ninguna: es la posicion del autor sobre lo humano.

\section{Método}
Un lugar vacio declarado como decision, con su razon escrita.

\section{Límites, y qué lo tumba}
Una definicion de lo humano puesta desde una sola mirada.

\section{Lo que queda abierto}
\textquestiondown{}Si no se puede definir lo humano, como se dice algo sobre el sin obligar a ver como uno?

\section*{Cómo se leyó esta tesis}
Esto se escribió en el repositorio man-made-empathy, del que este candado es una pieza: candados/02 el alma y el libre albedrio/083 lo humano se deja en blanco a proposito.md. La máquina que lo compuso lo hizo desde este mismo archivo, y no agregó nada que no estuviera escrito. Si algo de esto es falso, se dice en el mismo sitio, y se corrige con fecha.

\end{document}
